\subsection{When does the opportunity appear or disappear?}
The opportunity is a supported whole view ahead of a newly installed
selected frontier. In the headline RevTrace arm, 18 of 368 updates encounter
and retain one; computation-only encounters 17 and clears all 17. Under
synchronized arrivals, RevTrace encounters only 3 of 368 and covers 0.790~s
less than computation-only. A fresh staggered repeat on cohort~A also
changes the timing ranking: 226.320 versus 230.565~s, $I=-2.245$~s.
Compatible relative progress determines the benefit; the three positive
primary contrasts do not establish an invariant speed advantage.

\begin{table}[t]
\centering
\caption{Operating boundaries remain visible. Each workload has four
arms and three fresh processes (48 runs total). Pass counts cover all 12
runs; D99 ranges span three RevTrace observations, and timely-loss windows
list computation-only/RevTrace counts. Source mix expands with sessions. Full 48-point
plots and all endpoints remain in the artifact.}
\label{tab:operating_workloads}
\small\setlength{\tabcolsep}{3pt}
\begin{tabular}{@{}lccc@{}}
\toprule
Workload & D2/C8 pass & \shortstack{RevTrace\\D99 (s)} &
\shortstack{Timely-loss\\windows}\\
\midrule
8 sync &12/12 &1.426--1.528 &1,0,0 /0,0,0\\
8 staggered &12/12 &1.704--1.955 &1,2,1 /0,0,0\\
12 staggered &0/12 &2.197--2.299 &1,1,3 /0,0,0\\
16 staggered &0/12 &2.801--3.223 &10,12,9 /0,0,0\\
\bottomrule
\end{tabular}
\end{table}

Table~\ref{tab:operating_workloads} summarizes 36 predeclared development
runs (8 synchronized, 8 staggered, 16 staggered sessions) and 12 exploratory
midpoint runs (12 staggered). Each arm offers 121.205/185.300/250.180~s
at 8/12/16 sessions. These are explicit workloads with different source
mixes. At 8 sessions all arms meet whole-run D2/C8; RevTrace synchronized
display P99 is 1.426--1.528~s and staggered reaches 1.955~s. Two 0.215-s
tails remain uncredited despite their P99 passes. At 12 and 16 sessions
all runs fail D2 and pass C8. At 16, RevTrace's extra coverage over
computation-only is +18.910,+5.055,$-7.225$~s; P99 improves in one repeat.
Its continuity distinction is stable: compatible withdrawals are zero
versus 16/17/12.

A secondary analysis follows intervals already delivered within D2 whose
eventual certified coverage also meets C8. At each selected application,
the complete dependency path classifies compatible loss separately from
necessary revision. Loss ends when actual coverage returns or is
right-censored at the next selected root. Computation-only has 10/12/9
such windows at 16 sessions, with 24.903/23.635/15.736 audio-s$\times$wall-s
deficit; RevTrace has none. Holes intersect the original D2 read point
for 6.000/10.675/4.000~s of distinct audio. The 12-session strong witness
is not established: computation-only's smaller losses never include a
timely complete non-EOS interval under its frozen criterion.
\par\noindent\textbf{Finding 3: continuity and first delivery have different boundaries.}
Retaining a supported view removes an avoidable withdrawal; whether
enough work arrives before D2 remains a separate scheduling and load outcome.

\paragraph{Quality and cost}
Headline Qwen arms share WER 26/1908 (1.36\%) and selected histories.
Retrospective same-frontier whole-text agreement leaves their credits
unchanged. Elsewhere RevTrace's raw comparison removes 2~s in the 16-session third
repeat and 2~s in each midpoint's first two repeats, although the frozen
normalizer agrees. Whisper WER is 83.12\% at $\rho=1$ versus 11.53\% at
$\rho=3$; changing its feedback contract changes quality. Headline board
energy spans 59.16--59.73~J/audio-s across five arm medians. At 16 sessions, RevTrace reduces total
model time by 0.66--1.59~s over recomputation, while wall-time and energy
differences change sign. Retained progress is the measured benefit;
inference cost and overload are independent boundaries.
